\documentclass[../../../../master.tex]{subfiles}
\begin{document}

\monographleaf{Domain 06: Reinforcement Learning \& Alignment}{Policy Gradients, Model-Based RL, MARL \& LLM Alignment}{Policy Gradients, Actor-Critic, Model-Based RL \& MARL}{Dyna, MCTS, World Models, Offline RL \& Multi-Agent Reinforcement Learning}{\bookcite{reinforcement_learning_introduction_richard_s_su}{Ch.~8}, \bookcite{reinforement_learning_kevin_murphy}{Ch.~4--5 \& Sec.~7.5--7.7}, \bookcite{probabilisticml_kevinmurphy_advancedtopics}{Sec.~35.4--35.5}}

\section{Problem Formulation \& Unified Notation}
\label{sec:mbrl:formulation}

Model-free reinforcement learning (RL) estimates value functions or policies directly from environment interactions, often requiring millions of samples \newline \bookcite{probabilisticml_kevinmurphy_advancedtopics}{Sec.~35.4, p.~1164}. Model-based RL (MBRL) improves sample efficiency by learning an internal approximation of the environment's transition dynamics and reward function, enabling \emph{planning}---any computational process that takes a model as input and produces or improves a policy \newline \bookcite{reinforcement_learning_introduction_richard_s_su}{Ch.~8, p.~160}.

Consider a Markov Decision Process (MDP) $\mathcal{M} = (\States, \Actions, \Trans, R, \gamma, \mu_0)$ with state space $\States$, action space $\Actions$, transition kernel $\Trans(s' \mid s, a)$, reward function $R(s, a, s')$, discount factor $\gamma \in [0, 1)$, and initial distribution $\mu_0$. A model of the environment can be either a \emph{distribution model}, which outputs the full conditional probability $p(s', r \mid s, a)$, or a \emph{sample model}, which generates a single draw $(\hat{s}', \hat{r}) \sim \hat{\Trans}(\cdot \mid s, a)$ \bookcite{reinforcement_learning_introduction_richard_s_su}{Ch.~8, pp.~159--160}. Planning splits into two regimes:
\begin{itemize}
  \item \textbf{Background (offline) planning} (e.g., Dyna, Dreamer): uses simulated experience from a learned model $\hat{M}$ to update a global policy $\pi_\theta$ or value function $Q_w$ in the background \bookcite{reinforement_learning_kevin_murphy}{Sec.~4.3, p.~93}.
  \item \textbf{Decision-time (online) planning} (e.g., MCTS, MPC): waits until state $s_t$ is observed and searches forward over a finite horizon $H$ to select the immediate action $a_t$ \newline \bookcite{reinforcement_learning_introduction_richard_s_su}{Ch.~8, p.~180}, \newline \bookcite{probabilisticml_kevinmurphy_advancedtopics}{Sec.~35.4.1, p.~1164}.
\end{itemize}

When real-time interaction is prohibited, \textbf{offline RL} learns solely from a static dataset $\mathcal{D} = \{(s_i, a_i, r_i, s'_i)\}_{i=1}^m$ collected by a behavior policy $\pi_b$ \bookcite{reinforement_learning_kevin_murphy}{Sec.~7.7, p.~192}, \bookcite{probabilisticml_kevinmurphy_advancedtopics}{Sec.~35.5, p.~1171}. When $n \ge 2$ agents interact simultaneously, the MDP generalizes to a \textbf{Partially Observable Stochastic Game} (POSG) $(\mathcal{I}, \States, \{\Actions_i\}_{i=1}^n, \Trans, \{R_i\}_{i=1}^n, \{\mathcal{O}_i\}_{i=1}^n, \gamma)$ over joint actions $\vecb{a} = (a_1, \dots, a_n)$ and agent-specific observation histories $h_t^i = \sigma_i(\hat{h}_t)$ \bookcite{reinforement_learning_kevin_murphy}{Sec.~5.1--5.2, pp.~121, 127}.

\begin{crossbook}[Cross-Book Notation \& Pedagogical Reconciliation]
\label{tab:mbrl:notation}
\begin{tabularx}{\textwidth}{@{}l >{\raggedright\arraybackslash}X >{\raggedright\arraybackslash}X >{\raggedright\arraybackslash}X@{}}
\toprule
\textbf{Concept} & \textbf{Sutton \& Barto (Ch.~8)} & \textbf{Murphy RL (Ch.~4--5, 7)} & \textbf{Murphy Adv. (Sec.~35.4--35.5)} \\
\midrule
Dynamics / Model & $p(s', r \mid s, a)$ (distribution) vs. sample model & $T(s' \mid s, a)$, $M_w(z_t, a_t)$ (latent) & $p(s' \mid s, a)$, $R(s, a)$, $M_\theta$ \\
Latent World State & Tabular state $s \in \States$ & Deterministic $h_t$, stochastic $z_t$ (Dreamer); $z_t = e_\phi(o_t)$ (MuZero) & Latent $z_t$ (VAE/RSSM), deterministic $h_t$ \\
Behavior Policy & $b(a \mid s)$ (off-policy target $\pi$) & $\pi_b(a \mid s)$ or $\pi_\beta$ (offline dataset $\mathcal{D}$) & $\pi_b$ (behavior) vs. $\pi_e$ (evaluation) \\
Planning Horizon & $n$-step lookahead / search depth & Receding horizon $H$, search budget $m$ & Planning horizon $H$, terminal $\hat{V}(s_{t+H})$ \\
Multi-Agent History & Single-agent MDP (two-player zero-sum in MCTS) & Full history $\hat{h}_t$, agent history $h_t^i = \sigma_i(\hat{h}_t)$, opponents $\pi_{-i}$ & Single-agent focus; off-policy flow $d(s,a)$ \\
\bottomrule
\end{tabularx}
\end{crossbook}

\section{Core Mathematical Derivations \& Theorems}
\label{sec:mbrl:derivations}

\subsection{Dyna, Probabilistic Ensembles (PETS), and Compounding Error}
\label{sec:mbrl:dyna_pets}

\begin{coremath}[Dyna Integration, MPC Objective, and Compounding Model Error]
In tabular \textbf{Dyna-Q} \bookcite{reinforcement_learning_introduction_richard_s_su}{Ch.~8, pp.~161--164}, \bookcite{reinforement_learning_kevin_murphy}{Sec.~4.3.2, p.~96}, each real transition $(s, a, r, s')$ triggers three updates: (i) \emph{direct RL} on $Q(s, a)$, (ii) \emph{model learning} $\hat{M}(s, a) \leftarrow (s', r)$, and (iii) $N$ \emph{planning updates} on simulated transitions $(\tilde{s}', \tilde{r}) \sim \hat{M}(\tilde{s}, \tilde{a})$ sampled from previously visited pairs (or prioritized by Bellman residual $|\delta|$ in \emph{prioritized sweeping} \bookcite{reinforcement_learning_introduction_richard_s_su}{Ch.~8, p.~168}).

In continuous control with \textbf{Model Predictive Control (MPC)}, the agent solves at each step $t$ \bookcite{probabilisticml_kevinmurphy_advancedtopics}{Eq.~35.46, p.~1164}:
\begin{equation}
\label{eq:mbrl:mpc}
a^*_{t:t+H-1} = \argmax_{a_{t:t+H-1}} \E\!\left[\sum_{h=0}^{H-1} R(s_{t+h}, a_{t+h}) + \hat{V}(s_{t+H})\right],
\end{equation}
and executes only $a^*_t$. To prevent policy exploitation of overconfident point-estimate dynamics, \textbf{PETS} (Probabilistic Ensembles with Trajectory Sampling) \bookcite{probabilisticml_kevinmurphy_advancedtopics}{Sec.~35.4.4, p.~1168} models $p_\theta(s' \mid s, a)$ via an ensemble of $B$ probabilistic neural networks $\{\Normal(\mu_{\theta_b}(s, a), \Sigma_{\theta_b}(s, a))\}_{b=1}^B$, separating \emph{aleatoric} uncertainty (intra-member variance $\Sigma_{\theta_b}$) from \emph{epistemic} uncertainty (inter-member disagreement in $\mu_{\theta_b}$).

When a learned model has bounded single-step total variation error $\sup_{s,a} D_{\mathrm{TV}}(\hat{\Trans}(\cdot \mid s,a) \,\|\, \Trans(\cdot \mid s,a)) \le \epsilon_m$ and the policy diverges from the data-collection policy by $\epsilon_\pi$, the return error scales quadratically with horizon:
\begin{equation}
\label{eq:mbrl:compounding}
\abs{J_{\hat{\mathcal{M}}}(\pi) - J_{\mathcal{M}}(\pi)} \le \frac{2 R_{\max}}{(1-\gamma)^2}\left[\epsilon_\pi + \frac{\gamma}{1-\gamma}\epsilon_m\right], \qquad
\mathcal{E}_{\text{finite}}(H) = \mathcal{O}(H^2 \epsilon_m).
\end{equation}
\end{coremath}

\begin{intuition}[Why Truncated Rollouts (MBPO) Beat Full-Horizon Imagination]
Because single-step model errors compound quadratically as $\mathcal{O}(H^2 \epsilon_m)$ across long synthetic trajectories, rolling out an imperfect dynamics model from $s_0$ to termination $T$ injects severe bias. Branching short $k$-step model rollouts (e.g., $k \in [1, 5]$) from real states stored in the replay buffer $\mathcal{D}$---the Dyna/MBPO principle---caps model error at $\mathcal{O}(k^2 \epsilon_m)$ while still multiplying effective sample efficiency by $k$ \bookcite{reinforement_learning_kevin_murphy}{Sec.~4.3.2, p.~96}.
\end{intuition}

\subsection{Decision-Time Planning: MCTS, UCT, AlphaZero PUCT, and MuZero}
\label{sec:mbrl:mcts}

Monte Carlo Tree Search (MCTS) builds an asymmetric search tree rooted at $s_t$ through $K$ simulations of four phases: \emph{Selection}, \emph{Expansion}, \emph{Simulation/Evaluation}, and \emph{Backup} \newline \bookcite{reinforcement_learning_introduction_richard_s_su}{Ch.~8, pp.~185--187}, \newline \bookcite{reinforement_learning_kevin_murphy}{Sec.~4.2.2, pp.~86--88}.

\begin{coremath}[UCT, AlphaZero PUCT, and MuZero Latent Search]
At internal tree nodes, \textbf{UCT} (Upper Confidence Bounds for Trees) selects actions via \bookcite{reinforement_learning_kevin_murphy}{Eq.~4.1, p.~86}:
\begin{equation}
\label{eq:mbrl:uct}
a_t = \argmax_{a \in \Actions(s)} \left[ Q(s, a) + c \sqrt{\frac{\ln N(s)}{N(s, a)}} \right], \qquad N(s) = \sum_{b} N(s, b).
\end{equation}
\textbf{AlphaZero} replaces the uniform UCB bonus with a learned neural policy prior $P(s, a) = \pi_\theta(a \mid s)$ via the \textbf{PUCT} rule and eliminates Monte Carlo rollouts by evaluating leaf nodes with a value head $V_\theta(s)$ \bookcite{reinforement_learning_kevin_murphy}{Eqs.~4.2--4.4, pp.~86--88}:
\begin{equation}
\label{eq:mbrl:puct}
a_t = \argmax_{a \in \Actions(s)} \left[ Q(s, a) + c_{\text{puct}}\, P(s, a) \frac{\sqrt{N(s)}}{1 + N(s, a)} \right], \qquad
Q(s, a) \leftarrow Q(s, a) + \frac{u - Q(s, a)}{N(s, a)}.
\end{equation}
After $K$ simulations, the root visit distribution $\pi^{\text{MCTS}}_s(a) \propto N(s, a)^{1/\tau}$ serves as the policy target, and the network $(V_\theta, \pi_\theta)$ minimizes:
\begin{equation}
\label{eq:mbrl:alphazero_loss}
\mathcal{L}_{\text{AZ}}(\theta) = \E_{(s, \pi^{\text{MCTS}}_s, V^{\text{MCTS}}_s) \sim \mathcal{D}} \left[ \left(V^{\text{MCTS}}(s) - V_\theta(s)\right)^2 - \sum_{a} \pi^{\text{MCTS}}_s(a) \log \pi_\theta(a \mid s) \right].
\end{equation}
\textbf{MuZero} removes the requirement of a known simulator by learning an encoder $z_t = e_\phi(o_t)$ and latent dynamics/reward model $(z_{t+1}, \hat{r}_t) = M_w(z_t, a_t)$, jointly optimizing $(\theta, w, \phi)$ to match observed rewards $r_t$ alongside MCTS value and policy targets without reconstructing pixels \bookcite{reinforement_learning_kevin_murphy}{Eqs.~4.5--4.6, p.~88}:
\begin{equation}
\label{eq:mbrl:muzero_loss}
\mathcal{L}_{\text{MZ}}(\theta, w, \phi) = \mathcal{L}_{\text{AZ}}(\theta, w, \phi) + \E_{\mathcal{D}}\!\left[\left(r_t - \hat{r}_t(e_\phi(o_t), a_t)\right)^2\right].
\end{equation}
\end{coremath}

\subsection{Latent World Models (Dreamer RSSM) \& Successor Features}
\label{sec:mbrl:dreamer_sf}

Rather than predicting high-dimensional pixels directly or planning only at decision time, \textbf{Dreamer} learns a Recurrent State-Space Model (RSSM) combining a deterministic recurrent state $h_t = U(h_{t-1}, a_{t-1}, z_{t-1})$ with a stochastic latent $z_t$ \bookcite{reinforement_learning_kevin_murphy}{Sec.~4.4.1.3, pp.~98--100}, \bookcite{probabilisticml_kevinmurphy_advancedtopics}{Sec.~35.4.5, p.~1170}.

\begin{coremath}[Dreamer RSSM ELBO, Imagination Actor-Critic, and Successor Features]
Given conditional prior $\hat{z}_t \sim P(\hat{z}_t \mid h_t)$, posterior encoder $z_t \sim E(z_t \mid h_t, o_t)$, decoder $\hat{o}_t \sim D(\hat{o}_t \mid h_t, z_t)$, and reward head $\hat{r}_t \sim R(\hat{r}_t \mid h_t, z_t)$, the world model minimizes the variational ELBO loss \bookcite{reinforement_learning_kevin_murphy}{Eqs.~4.21--4.24, p.~100}:
\begin{equation}
\label{eq:mbrl:dreamer_wm}
\mathcal{L}_{\text{WM}} = \E_{q(z_{1:T} \mid h_0, o_{1:T}, a_{1:T})}\!\left[\sum_{t=1}^T \left( -\beta_o \ln D(o_t \mid z_t, h_t) + \beta_z \KL{E(z_t \mid h_t, o_t)}{P(\hat{z}_t \mid h_t)} \right)\right].
\end{equation}
Inside latent imagination (unrolling $h_{t+1} = U(h_t, a_t, \hat{z}_t)$ without decoding pixels), the critic and actor optimize GAE $\lambda$-returns $G_t^\lambda = r_t + \gamma\left((1-\lambda)V(h_{t+1}) + \lambda G_{t+1}^\lambda\right)$ \bookcite{reinforement_learning_kevin_murphy}{Eqs.~4.25--4.27, p.~100}:
\begin{equation}
\label{eq:mbrl:dreamer_ac}
\mathcal{L}_{\text{critic}} = \sum_{t=1}^T \left(V(h_t) - \mathrm{sg}(G_t^\lambda)\right)^2, \qquad
\mathcal{L}_{\text{actor}} = -\sum_{t=1}^T \mathrm{sg}\!\left(G_t^\lambda - V(h_t)\right)\log \pi(a_t \mid h_t).
\end{equation}
For fast transfer across tasks with linear rewards $R(s, \vecb{w}) = \vecb{\phi}(s)\T \vecb{w}$, \textbf{Successor Features (SFs)} decouple dynamics from reward via the discounted feature occupancy \bookcite{reinforement_learning_kevin_murphy}{Eqs.~4.60--4.71, pp.~115--116}:
\begin{equation}
\label{eq:mbrl:sf_gpi}
\vecb{\psi}^\pi(s, a) = \E_\pi\!\left[\sum_{t=0}^\infty \gamma^t \vecb{\phi}(s_{t+1}) \given s_0=s, a_0=a\right], \qquad
Q^{\pi, \vecb{w}}(s, a) = \vecb{\psi}^\pi(s, a)\T \vecb{w}.
\end{equation}
Given a library of policies $\{\pi_i\}_{i=1}^N$ and a new task $\vecb{w}_{\text{new}}$, \textbf{Generalized Policy Improvement (GPI)} selects $a^*(s; \vecb{w}_{\text{new}}) = \argmax_a \max_i \vecb{\psi}^{\pi_i}(s, a)\T \vecb{w}_{\text{new}}$, guaranteeing $Q^{\pi}(s, a) \ge \max_i Q^{\pi_i}(s, a)$ \bookcite{reinforement_learning_kevin_murphy}{Sec.~4.5.3.1, p.~116}.
\end{coremath}

\subsection{Offline RL Conservatism \& Off-Policy Importance Corrections}
\label{sec:mbrl:offline}

In offline RL, standard Q-learning suffers from \emph{extrapolation error}: $\max_{a'} Q(s', a')$ queries out-of-distribution (OOD) actions whose overestimation bias cannot be corrected by online exploration \bookcite{reinforement_learning_kevin_murphy}{Sec.~7.7.1, p.~193}.

\begin{coremath}[CQL, Policy Constraints, Offline MBRL, and DualDICE]
\textbf{Conservative Q-Learning (CQL)} \bookcite{reinforement_learning_kevin_murphy}{Eqs.~7.56--7.59, p.~195} augments the Bellman error $\mathcal{E}(B, w)$ with a minimax penalty that pushes down $Q$-values under a high-value proposal $\mu(a \mid s)$ while pushing up $Q$-values on dataset actions $\pi_b(a \mid s)$:
\begin{equation}
\label{eq:mbrl:cql}
\mathcal{L}_{\text{CQL}}(w) = \alpha \left(\E_{s \sim \mathcal{D},\, a \sim \mu(\cdot \mid s)}[Q_w(s, a)] - \E_{(s,a) \sim \mathcal{D}}[Q_w(s, a)] + \mathcal{R}(\mu)\right) + \frac{1}{2}\mathcal{E}(B, w).
\end{equation}
Setting $\mathcal{R}(\mu) = -\KL{\mu}{\rho}$ with uniform prior $\rho$ yields the soft log-sum-exp penalty $\log \sum_a \exp(Q_w(s, a))$. Alternatively:
\begin{itemize}
  \item \textbf{Implicit KL Constraints (AWR/AWAC)}: solve $\pi_{k+1}(a \mid s) \propto \pi_b(a \mid s) \exp(\frac{1}{\alpha} Q^{\pi_k}(s, a))$ and project via advantage-weighted regression \bookcite{reinforement_learning_kevin_murphy}{Eqs.~7.51--7.52, p.~194}.
  \item \textbf{Offline MBRL (MOPO/MOREL)}: penalize rewards by model uncertainty $u(s, a)$, $\tilde{R}(s, a) = R(s, a) - \lambda u(s, a)$, yielding a guaranteed lower bound on true return when $D(M_\theta, M^*) \le u(s, a)$ \bookcite{reinforement_learning_kevin_murphy}{Sec.~7.7.2, p.~195}.
  \item \textbf{Sequence Modeling (Decision Transformer)}: autoregressively predicts actions conditioned on return-to-go $\mathrm{RTG}_t = \sum_{k=t}^T r_k$ via $\argmax_\theta \E_{\mathcal{D}}[\log \pi_\theta(a_t \mid s_{0:t}, a_{0:t-1}, \mathrm{RTG}_{0:t})]$ \bookcite{reinforement_learning_kevin_murphy}{Eq.~7.60, p.~196}.
\end{itemize}
For off-policy evaluation without exponential trajectory importance products $\prod_t \frac{\pi_e(a_t \mid s_t)}{\pi_b(a_t \mid s_t)}$ (the \emph{curse of horizon}), \textbf{DualDICE} estimates the marginalized stationary density ratio $\zeta^*(s, a) = p_\pi^\infty(s, a) / p_{\mathcal{D}}(s, a)$ via a convex-concave saddle point \bookcite{probabilisticml_kevinmurphy_advancedtopics}{Eqs.~35.56--35.58, p.~1175}:
\begin{equation}
\label{eq:mbrl:dualdice}
J(\pi) = \frac{1}{1-\gamma}\E_{(s,a)\sim p_{\mathcal{D}}}[\zeta^*(s, a) R(s, a)].
\end{equation}
\end{coremath}

\subsection{Multi-Agent RL: Equilibria, CTDE (VDN/QMIX/MAPPO), and Regret}
\label{sec:mbrl:marl}

In a stochastic game, agent $i$'s expected return $U_i(\pi_i, \pi_{-i}) = \E_{\hat{h} \sim p(\cdot \mid \pi)}[\sum_{t=0}^\infty \gamma^t R_i(s_t, \vecb{a}_t, s_{t+1})]$ depends on all opponents' policies $\pi_{-i}$ \bookcite{reinforement_learning_kevin_murphy}{Eqs.~5.11--5.13, p.~127}. A joint policy $\pi^*$ is a \textbf{Nash Equilibrium (NE)} iff no agent can unilaterally improve: $\forall i, \pi'_i,\, U_i(\pi'_i, \pi^*_{-i}) \le U_i(\pi^*)$ \bookcite{reinforement_learning_kevin_murphy}{Eq.~5.22, p.~128}. In two-player zero-sum games, closeness to minimax NE is measured by \emph{exploitability}: $\frac{1}{2}[\max_{\pi'_1} J(\pi'_1, \pi_2) - \min_{\pi'_2} J(\pi_1, \pi'_2)]$ \bookcite{reinforement_learning_kevin_murphy}{Eq.~5.21, p.~128}.

\begin{coremath}[CTDE Value Decomposition (VDN/QMIX), Multi-Agent Actor-Critic, and MMDS]
Under \textbf{Centralized Training with Decentralized Execution (CTDE)}, agents condition execution on local histories $h_t^i$ while training critics with global state/central information $c_t$. For cooperative common-reward games, value factorization requires the \textbf{Individual-Global-Max (IGM)} condition \bookcite{reinforement_learning_kevin_murphy}{Eq.~5.45, p.~136}:
\begin{equation}
\label{eq:mbrl:igm}
\argmax_{\vecb{a}} Q_{\text{tot}}(\vecb{h}_t, c_t, \vecb{a}; \theta) = \left(\argmax_{a^1} Q_1(h_t^1, a^1; \theta_1), \dots, \argmax_{a^n} Q_n(h_t^n, a^n; \theta_n)\right).
\end{equation}
\textbf{VDN} enforces IGM via linear additivity $Q_{\text{tot}} = \sum_{i=1}^n Q_i(h_t^i, a_t^i; \theta_i)$, whereas \textbf{QMIX} uses a hypernetwork-parameterized monotonic mixing network $Q_{\text{tot}} = f_{\text{mix}}(Q_1, \dots, Q_n; c_t)$ with non-negative weights $\frac{\partial f_{\text{mix}}}{\partial Q_i} \ge 0$ \bookcite{reinforement_learning_kevin_murphy}{Eqs.~5.46--5.49, p.~137}.

In policy-gradient MARL (\textbf{MAA2C / MAPPO}), the multi-agent policy gradient uses centralized critics $Q_i^\pi(\hat{h}, \vecb{a})$ or $V_i^\pi(\hat{h})$ to reduce variance \bookcite{reinforement_learning_kevin_murphy}{Eqs.~5.39--5.40, p.~134}:
\begin{equation}
\label{eq:mbrl:mappo}
\nabla_{\theta_i} J(\theta_{1:n}) \propto \E_{\hat{h} \sim p(\hat{h} \mid \pi),\, \vecb{a} \sim \pi}\!\left[\left(Q_i^\pi(\hat{h}, \vecb{a}) - V_i^\pi(\hat{h})\right) \nabla_{\theta_i} \log \pi_i(a_i \mid h_i; \theta_i)\right].
\end{equation}
For imperfect-information games, \textbf{Magnetic Mirror Descent Search (MMDS)} and Counterfactual Regret Minimization stabilize self-play toward NE via regularized local policy updates toward a magnet $\rho$ \bookcite{reinforement_learning_kevin_murphy}{Eq.~5.70, p.~143}:
\begin{equation}
\label{eq:mbrl:mmds}
\pi_{\text{new}} = \argmax_{\pi' \in \Delta(\Actions)} \left[ \inner{\pi'}{q} - \alpha \KL{\pi'}{\rho} - \frac{1}{\eta}\KL{\pi'}{\pi} \right].
\end{equation}
\end{coremath}

Beyond flat single-step MDPs, \textbf{Hierarchical RL (Options)} models temporally extended macro-actions $\omega = (\mathcal{I}_\omega, \pi_\omega, \beta_\omega)$ over semi-MDPs \bookcite{reinforement_learning_kevin_murphy}{Sec.~7.5.1, p.~184}, while \textbf{MaxEnt Inverse RL} recovers reward functions $R_\theta$ from expert demonstrations $\mathcal{D}_{\text{exp}}$ via energy-based trajectory likelihood $p(\tau) \propto \exp(\sum_{t=0}^{T-1} R_\theta(s_t, a_t))$ \newline \bookcite{reinforement_learning_kevin_murphy}{Eqs.~7.44--7.45, p.~191}.

\section{Algorithmic Implementation \& Comparative Synthesis}
\label{sec:mbrl:algorithms}

Algorithm~\ref{alg:mbrl:unified} unifies decision-time MCTS (AlphaZero/MuZero) with background Dyna/MBPO rollouts.

\begin{algorithm}[htbp]
\caption{Unified AlphaZero/MuZero MCTS \& Dyna-MBPO Planning Loop}
\label{alg:mbrl:unified}
\KwIn{Initial state/obs $o_0$, representation $e_\phi$, dynamics/reward model $M_w$, policy/value $(P_\theta, V_\theta)$, MCTS budget $K$, Dyna rollouts $N$, horizon $k_{\text{branch}}$}
\For{$t = 0, 1, 2, \dots$}{
  Encode root latent $z_t \leftarrow e_\phi(o_t)$ \tcp*{Identity $e_\phi(o_t)=s_t$ in AlphaZero}
  \For{$m = 1, \dots, K$ \tcp*{Decision-time MCTS}}{
    Initialize path $\tau \leftarrow [\,]$, $z \leftarrow z_t$\;
    \While{$z$ is expanded in tree}{
      Select $a \leftarrow \argmax_{a'} \left[Q(z, a') + c_{\text{puct}} P_\theta(a' \mid z)\frac{\sqrt{N(z)}}{1 + N(z, a')}\right]$\;
      Simulate $(z', \hat{r}) \leftarrow M_w(z, a)$; append $(z, a, \hat{r})$ to $\tau$; $z \leftarrow z'$\;
    }
    Expand leaf $z$ with prior $P_\theta(\cdot \mid z)$ and evaluate $u \leftarrow V_\theta(z)$\;
    \For{$(z_j, a_j, \hat{r}_j) \in \mathrm{reverse}(\tau)$ \tcp*{Backup}}{
      $u \leftarrow \hat{r}_j + \gamma u$; \quad $N(z_j, a_j) \leftarrow N(z_j, a_j) + 1$\;
      $Q(z_j, a_j) \leftarrow Q(z_j, a_j) + \frac{u - Q(z_j, a_j)}{N(z_j, a_j)}$\;
    }
  }
  Sample $a_t \sim \pi^{\text{MCTS}}_{z_t}(a) \propto N(z_t, a)^{1/\tau}$; execute $a_t$, observe $(r_t, o_{t+1})$, store in $\mathcal{D}$\;
  \For{$n = 1, \dots, N$ \tcp*{Background Dyna / MBPO truncated rollouts}}{
    Sample real state $o \sim \mathcal{D}$, encode $\tilde{z}_0 = e_\phi(o)$, unroll $k_{\text{branch}}$ steps via $M_w$ and $\pi_\theta$\;
    Update $(\theta, w, \phi)$ via Eq.~\eqref{eq:mbrl:muzero_loss} or Dreamer actor-critic Eqs.~\eqref{eq:mbrl:dreamer_wm}--\eqref{eq:mbrl:dreamer_ac}\;
  }
}
\end{algorithm}

\begin{table}[htbp]
\centering
\small
\caption{Comparative synthesis across model-free, Dyna/MBPO, MCTS, latent world models, and offline RL.}
\label{tab:mbrl:comparison}
\begin{tabularx}{\textwidth}{@{}l >{\raggedright\arraybackslash}p{2.2cm} l >{\raggedright\arraybackslash}p{2.0cm} >{\raggedright\arraybackslash}X@{}}
\toprule
\textbf{Paradigm} & \textbf{Planning Mode} & \textbf{Sample Eff.} & \textbf{Compute} & \textbf{Primary Bottleneck} \\
\midrule
Model-Free (DQN/PPO) & None & Low ($\mathcal{O}(10^7)$) & $\mathcal{O}(1)$ forward & Sample complexity \\
Dyna-Q / MBPO & Background ($k$-step) & High & $\mathcal{O}(1)$ forward & Compounding $\mathcal{O}(k^2\epsilon_m)$ \\
MCTS (AlphaZero/MuZero) & Decision-time (tree) & High (self-play) & $\mathcal{O}(K \cdot d)$ search & Discrete/sampled branching \\
Latent WM (Dreamer) & Background (latent) & Very High & $\mathcal{O}(1)$ policy & Obs. recon. / KL collapse \\
Offline RL (CQL/MOPO) & Static dataset $\mathcal{D}$ & Fixed ($0$ online) & $\mathcal{O}(1)$ forward & OOD extrapolation \\
\bottomrule
\end{tabularx}
\end{table}

\begin{crossbook}[Cross-Book Synthesis: Planning, Representation, and Distribution Shift]
Sutton \& Barto \bookcite{reinforcement_learning_introduction_richard_s_su}{Ch.~8, pp.~161, 188} frame model-based and model-free RL as endpoints of a single spectrum: both compute backed-up value targets, differing only in whether experience is real or simulated. Murphy \bookcite{reinforement_learning_kevin_murphy}{Sec.~4.4.2, p.~102} highlights the \emph{objective mismatch problem} of pixel-reconstruction world models---maximizing observation likelihood wastes capacity on task-irrelevant distractors---motivating value-equivalent latent models like MuZero \bookcite{reinforement_learning_kevin_murphy}{Sec.~4.2.2.2, p.~88} and non-generative self-prediction losses. Finally, Murphy Advanced Topics \bookcite{probabilisticml_kevinmurphy_advancedtopics}{Sec.~35.4.6--35.5, pp.~1171--1175} unifies model error in MBRL with covariate shift in off-policy/offline RL: in both settings, policy improvement shifts the state-action visitation distribution away from the training support $\mathcal{D}$, requiring either epistemic uncertainty penalties (PETS/MOPO) or conservative density-ratio constraints (CQL/DualDICE).
\end{crossbook}

\section{Visual Diagnostics \& Geometric Structure}
\label{sec:mbrl:visuals}

\monographfig[0.64\textwidth]{reinforcement_learning_introduction_richard_s_su_p185_fig1.png}{The general Dyna architecture integrating real experience (direct RL and model learning) with simulated experience (planning updates) \bookcite{reinforcement_learning_introduction_richard_s_su}{Ch.~8, Fig.~8.1, p.~163}.}{fig:mbrl:dyna_sutton}

\monographfig[0.60\textwidth]{reinforement_learning_kevin_murphy_p99_fig1.png}{Dreamer RSSM world model factor graph (1-step unroll): deterministic state $h_t$, conditional prior $\hat{z}_t \sim P_0(\cdot \mid h_t)$, posterior encoder $z_t \sim E(\cdot \mid h_t, o_t)$, decoder $D$, and KL/reconstruction losses \bookcite{reinforement_learning_kevin_murphy}{Ch.~4, Fig.~4.5, p.~99}.}{fig:mbrl:dreamer_factor}

\monographfig[0.86\textwidth]{reinforement_learning_kevin_murphy_p192_fig1.png}{Comparison of (a) online on-policy RL, (b) online off-policy RL with replay buffer $\mathcal{D}$, and (c) offline RL with static dataset $\mathcal{D}$ collected by $\pi_\beta$ \bookcite{reinforement_learning_kevin_murphy}{Sec.~7.7, Fig.~7.6, p.~192}.}{fig:mbrl:offline_paradigms}

Figure~\ref{fig:mbrl:dyna_world_model} synthesizes the background Dyna/Dreamer latent imagination loop alongside the four-phase decision-time MCTS/MuZero search cycle, while Figure~\ref{fig:mbrl:compounding_error} quantifies the $\mathcal{O}(H^2 \epsilon_m)$ compounding error trade-off that motivates truncated rollouts.

\begin{figure}[htbp]
\centering
\resizebox{\textwidth}{!}{%
\begin{tikzpicture}[
  node distance=1.0cm and 1.1cm,
  >=Stealth,
  every node/.style={font=\footnotesize\sffamily}
]
  % Left: Dyna & Latent World Model Loop
  \node[monoBox, minimum width=2.5cm, text width=2.3cm, align=center] (env) {Real Env.\\$(s_t, a_t, r_t, s_{t+1})$};
  \node[monoTealBox, right=2.1cm of env, minimum width=2.7cm, text width=2.5cm, align=center] (wm) {World Model $\hat{M}_w$\\$\hat{P}(z' \mid z, a),\, \hat{R}$};
  \node[monoIndigoBox, above=1.4cm of wm, minimum width=2.7cm, text width=2.5cm, align=center] (imag) {Imagined Rollout\\$\tilde{\tau}_{1:H}$ (Dreamer)};
  \node[monoBox, above=1.4cm of env, minimum width=2.5cm, text width=2.3cm, align=center] (pol) {Policy \& Value\\$(\pi_\theta, V_\theta, Q_w)$};

  \draw[monoArrow] (env) -- node[below=2pt, align=center, font=\scriptsize] {Model Learning\\(ELBO / $\mathcal{L}_r$)} (wm);
  \draw[monoArrow] (wm) -- node[right=2pt, align=center, font=\scriptsize] {Latent\\Unroll} (imag);
  \draw[monoArrow] (imag) -- node[above=2pt, align=center, font=\scriptsize] {Background\\Planning} (pol);
  \draw[monoArrow] (env) -- node[left=2pt, align=center, font=\scriptsize] {Direct RL\\(Real $\mathcal{D}$)} (pol);
  \draw[monoArrow, dashed, color=monographTeal] (pol.south east) -- node[pos=0.48, sloped, above, font=\scriptsize] {Action $a_t$} (env.north east);

  % Right: 4-Phase MCTS / MuZero Loop
  \node[monoIndigoBox, right=1.6cm of imag, minimum width=2.2cm, text width=2.05cm, align=center] (sel) {\textbf{1. Select}\\PUCT (Eq.~\ref{eq:mbrl:puct})};
  \node[monoTealBox, right=0.9cm of sel, minimum width=2.2cm, text width=2.05cm, align=center] (exp) {\textbf{2. Expand}\\$z' = M_w(z, a)$};
  \node[monoBox, below=1.4cm of exp, minimum width=2.2cm, text width=2.05cm, align=center] (eval) {\textbf{3. Evaluate}\\$u = V_\theta(z')$};
  \node[monoWarnBox, below=1.4cm of sel, minimum width=2.2cm, text width=2.05cm, align=center] (back) {\textbf{4. Backup}\\Update $N, Q$};

  \draw[monoArrow] (sel) -- (exp);
  \draw[monoArrow] (exp) -- (eval);
  \draw[monoArrow] (eval) -- (back);
  \draw[monoArrow] (back) -- (sel);

  \draw[monoArrow, dashed, color=monographIndigo] (wm.east) -- node[below=2pt, font=\scriptsize, align=center] {Decision-Time\\Search} (back.west);
\end{tikzpicture}%
}
\caption{Unified architecture connecting background planning (Dyna / Dreamer RSSM latent imagination, left) with decision-time planning (4-phase MCTS / MuZero loop, right).}
\label{fig:mbrl:dyna_world_model}
\end{figure}

\begin{figure}[htbp]
\centering
\begin{tikzpicture}
\begin{axis}[
  width=0.84\textwidth,
  height=6.0cm,
  xlabel={Model Rollout Horizon $H$ (steps)},
  ylabel={Normalized Return Error / Objective Gap},
  xmin=1, xmax=25,
  ymin=0, ymax=1.2,
  grid=both,
  grid style={monographRule!40},
  legend pos=north west,
  legend style={font=\footnotesize, draw=monographRule, fill=white},
  tick label style={font=\footnotesize},
  label style={font=\small\bfseries, color=monographNavy}
]
  % Compounding model error O(H^2 * eps_m)
  \addplot[color=monographCrimson, very thick, domain=1:25, samples=50]
    {0.0018 * x^2};
  \addlegendentry{Compounding Model Error $\propto H^2 \epsilon_m$}

  % Value bootstrapping / sample variance reduction as H grows
  \addplot[color=monographTeal, very thick, dashed, domain=1:25, samples=50]
    {0.55 * exp(-0.22 * x) + 0.04};
  \addlegendentry{Value-Bias / Variance Reduction $\propto e^{-\alpha H}$}

  % Total MBPO / Dyna Effective Error (sum)
  \addplot[color=monographNavy, ultra thick, domain=1:25, samples=50]
    {0.0018 * x^2 + 0.55 * exp(-0.22 * x) + 0.04};
  \addlegendentry{Total Return Error (Optimal $H^* \approx 6$)}

  \draw[dashed, color=monographIndigo, thick] (axis cs:6.3, 0) -- (axis cs:6.3, 0.25);
  \draw[-{Stealth}, color=monographIndigo, thick] (axis cs:6.3, 0.44) -- (axis cs:6.3, 0.26);
  \node[anchor=south, font=\scriptsize\sffamily, color=monographIndigo, fill=white, inner sep=2pt] at (axis cs:6.3, 0.44) {Truncated MBPO Sweet Spot ($H^*$)};
\end{axis}
\end{tikzpicture}
\caption{Trade-off between compounding model error $\mathcal{O}(H^2 \epsilon_m)$ and value-function bias reduction across rollout horizon $H$, illustrating why truncated Dyna/MBPO branches minimize total return error.}
\label{fig:mbrl:compounding_error}
\end{figure}

\section{Numerical Pitfalls, Edge Cases \& Failure Modes}
\label{sec:mbrl:pitfalls}

\begin{pitfall}[The Deadly Triad \& Off-Policy Divergence (Baird's Counterexample)]
Combining \emph{function approximation}, \emph{bootstrapping}, and \emph{off-policy} updates---the \textbf{deadly triad}---can cause TD(0) and Q-learning weights to diverge to $\infty$ even on a 7-state deterministic MDP with zero rewards and $\gamma = 0.99$ \bookcite{probabilisticml_kevinmurphy_advancedtopics}{Sec.~35.5.3, Fig.~35.10, p.~1176}.
\textbf{Remedy:} Use target networks, gradient-TD / Emphatic-TD corrections, or behavior-regularized offline updates (CQL/IQL).
\end{pitfall}

\begin{pitfall}[Model Exploitation, Objective Mismatch \& Environment Non-Stationarity]
If a learned model $\hat{M}$ is inaccurate or the environment changes (e.g., blocking/shortcut maze tasks), greedy planning exploits phantom high-reward transitions \bookcite{reinforcement_learning_introduction_richard_s_su}{Ch.~8, pp.~166--168}. Similarly, generative pixel-reconstruction models suffer from \emph{objective mismatch} when capacity is spent on irrelevant distractors \bookcite{reinforement_learning_kevin_murphy}{Sec.~4.4.2.1, p.~102}.
\textbf{Remedy:} Incorporate exploration bonuses on state-action staleness (Dyna-Q+), uncertainty penalties via PETS ensembles or MOPO ($\tilde{R} = R - \lambda u$), or value-equivalent latent losses (MuZero).
\end{pitfall}

\begin{pitfall}[Offline Extrapolation Error \& Stochastic Return-Conditioning Bias]
Unconstrained off-policy Q-learning on a static dataset $\mathcal{D}$ overestimates OOD actions; conversely, return-conditioned sequence models (Decision Transformer) suffer from \emph{optimism bias} in stochastic MDPs, mistaking lucky high-return trajectories for optimal actions \bookcite{reinforement_learning_kevin_murphy}{Sec.~7.7.1 \& 7.7.3, pp.~193, 196}.
\textbf{Remedy:} Enforce pessimistic value penalties (CQL, Eq.~\eqref{eq:mbrl:cql}), support constraints (MMD/AWAC), or condition on learned expectiles / latent plans rather than raw RTG.
\end{pitfall}

\begin{pitfall}[MARL Non-Stationarity \& Independent Learning Oscillation]
In multi-agent environments, simultaneous policy updates render the effective transition dynamics non-stationary from each agent's local perspective, causing independent gradient ascent (IGA) or IPPO to cycle indefinitely (e.g., in Rock-Paper-Scissors) \bookcite{reinforement_learning_kevin_murphy}{Sec.~5.3.2.4, p.~135}.
\textbf{Remedy:} Use CTDE critics (QMIX/MAPPO), two-timescale learning rates (WoLF-PHC), or regularized self-play (MMDS/CFR, Eq.~\eqref{eq:mbrl:mmds}) with historical opponent pools.
\end{pitfall}

\end{document}
